% =====================================================
% 题型：解三角形面积与边长解答题
% =====================================================
% 难度：★★ (中档)
% =====================================================
\newcommand{\qTriangleAreaOneStem}{%
  在 \(\triangle ABC\) 中，角 \(A, B, C\) 的对边分别为 \(a, b, c\)，且 \(a \cos B = c + b \cos 2A\)。\\
  (1) 求 \(A\)；\\
  (2) 若 \(D\) 为 \(BC\) 的中点，\(AD = 3\)，\(\triangle ABC\) 的面积为 \(3\sqrt{3}\)，求 \(a\)。%
}

\newcommand{\qTriangleAreaOneAnswer}{%
  (1) \(A = \dfrac{\pi}{3}\)；(2) \(a = 2\sqrt{3}\)。%
}

\newcommand{\qTriangleAreaOneSolution}{%
  (1) 解：因为 $a \cos B = c + b \cos 2A$，
  
  由正弦定理 $\dfrac{a}{\sin A} = \dfrac{b}{\sin B} = \dfrac{c}{\sin C}$，可得
  \[ \sin A \cos B = \sin C + \sin B \cos 2A. \]
  
  因为在 $\triangle ABC$ 中，$\sin C = \sin(A+B) = \sin A \cos B + \cos A \sin B$，
  
  代入上式化简得
  \[ \cos A \sin B + \sin B \cos 2A = 0. \]
  
  即
  \[ \sin B (\cos A + \cos 2A) = 0. \]
  
  因为 $B \in (0, \pi)$，所以 $\sin B > 0$，
  
  从而有
  \[ \cos A + \cos 2A = 0. \]
  
  由二倍角公式 $\cos 2A = 2\cos^2 A - 1$，代入上式得
  \[ 2\cos^2 A + \cos A - 1 = 0, \]
  
  即
  \[ (2\cos A - 1)(\cos A + 1) = 0. \]
  
  因为 $A \in (0, \pi)$，所以 $\cos A \in (-1, 1)$，即 $\cos A + 1 \neq 0$，
  
  故
  \[ \cos A = \frac{1}{2}, \]
  
  解得
  \[ A = \frac{\pi}{3}. \]
  
  (2) 解：因为 $\triangle ABC$ 的面积为 $3\sqrt{3}$，
  
  所以
  \[ S = \frac{1}{2}bc\sin A = 3\sqrt{3}, \]
  
  即
  \[ \frac{1}{2}bc\sin\frac{\pi}{3} = \frac{\sqrt{3}}{4}bc = 3\sqrt{3}, \]
  
  解得
  \[ bc = 12. \]
  
  因为 $D$ 为 $BC$ 的中点，由向量加法可得 $\overrightarrow{AD} = \frac{1}{2}(\overrightarrow{AB} + \overrightarrow{AC})$，
  
  两边平方得
  \[ \overrightarrow{AD}^2 = \frac{1}{4}(\overrightarrow{AB}^2 + \overrightarrow{AC}^2 + 2\overrightarrow{AB}\cdot\overrightarrow{AC}), \]
  
  即
  \[ AD^2 = \frac{1}{4}(c^2 + b^2 + 2bc\cos A). \]
  
  代入 $AD = 3, \cos A = \frac{1}{2}$，得
  \[ 9 = \frac{1}{4}(b^2 + c^2 + bc), \]
  
  所以
  \[ b^2 + c^2 + bc = 36. \]
  
  将 $bc = 12$ 代入，得
  \[ b^2 + c^2 = 24. \]
  
  在 $\triangle ABC$ 中，由余弦定理得
  \[ a^2 = b^2 + c^2 - 2bc\cos A, \]
  
  代入数据得
  \[ a^2 = 24 - 2 \times 12 \times \frac{1}{2} = 12. \]
  
  因为 $a > 0$，所以
  \[ a = 2\sqrt{3}. \]
}

\newcommand{\qTriangleAreaOneDiagnosis}{%
  （留白待收集学生作答后补充）%
}

\newcommand{\qTriangleAreaOneTeachingNotes}{%
  精讲。第(1)问考察边角互化及和差角公式，重点提示 \(\sin C = \sin(A+B)\) 的展开；第(2)问考察中线长公式或向量平方求模，以及面积公式与余弦定理的结合使用。%
}
